\usepackage{geometry}
\IfFileExists{ukrainian.ldf}{%
    \usepackage[ukrainian]{babel}%
}{%
    \usepackage{polyglossia}\setmainlanguage{ukrainian}%
}
\usepackage{fontspec}
\usepackage{xcolor}
\usepackage{fancyhdr}
\usepackage{setspace}
\usepackage{array}
\usepackage{ragged2e}
\usepackage{titlesec}
\usepackage{indentfirst}
\usepackage[titles]{tocloft}
\usepackage{enumitem}
\usepackage{listings}
\usepackage{chngcntr}
\usepackage{xltabular}
\usepackage{graphicx}
\usepackage{caption}
\usepackage{color}
\usepackage{hyperref}

\usepackage{amsmath,amssymb}

% Діаграми (TikZ) вбудовано безпосередньо в документ: окремої компіляції assets/*.tex не потрібно
\usepackage{tikz}
\usetikzlibrary{arrows.meta, positioning, shapes.geometric, shapes.multipart, calc}
% Шрифти лежать у проєкті (каталог fonts/): Liberation Serif метрично сумісний із Times New Roman і містить кирилицю;
% так збірка не залежить від шрифтів, установлених у системі чи на Overleaf.
\setmainfont{LiberationSerif}[Path=./fonts/, Extension=.ttf, UprightFont=*-Regular, BoldFont=*-Bold, ItalicFont=*-Italic, BoldItalicFont=*-BoldItalic]

\hyphenpenalty=10000
\exhyphenpenalty=10000
\sloppy
% Моноширинний шрифт з підтримкою кирилиці (для лістингів та промтів у додатках)
\setmonofont{DejaVuSansMono}[Path=./fonts/, Extension=.ttf, UprightFont=*, BoldFont=*-Bold, Scale=0.85]
